% Created 2026-06-02 Tue 06:48
% Intended LaTeX compiler: pdflatex
\documentclass[11pt]{article}
\usepackage[utf8]{inputenc}
\usepackage[T1]{fontenc}
\usepackage{graphicx}
\usepackage{longtable}
\usepackage{wrapfig}
\usepackage{rotating}
\usepackage[normalem]{ulem}
\usepackage{amsmath}
\usepackage{amssymb}
\usepackage{capt-of}
\usepackage{hyperref}
\author{Arun Gopalan}
\date{\today}
\title{Mofography native granular visualization (main tab + PLD / percolation)}
\hypersetup{
 pdfauthor={Arun Gopalan},
 pdftitle={Mofography native granular visualization (main tab + PLD / percolation)},
 pdfkeywords={},
 pdfsubject={Runnable Org snippets mirroring mgr\textsubscript{cli}\textsubscript{py} interactive app — show\textsubscript{mof}\textsubscript{plotly}, show\textsubscript{mof}\textsubscript{threejs}, Figure3JS chain, show\textsubscript{percolation}\textsubscript{threejs}.},
 pdfcreator={Emacs 30.2 (Org mode 9.7.39)}, 
 pdflang={English}}
\begin{document}

\maketitle
\tableofcontents

\section{Overview}
\label{sec:orgeb213a1}

Reproduce the \textbf{\textbf{Visualization}} tab and \textbf{\textbf{PLD \& Percolation}} tab from \texttt{mgr\_cli\_py} using native \texttt{mofography} APIs only.

\begin{center}
\begin{tabular}{llll}
Layer & Plotly & Three.js (granular) & Three.js (one-liner)\\
\hline
Main MOF scene & \texttt{show\_mof\_plotly} & \texttt{Figure3JS} chain & \texttt{show\_mof\_threejs}\\
PLD / percolation & \texttt{visualize.plot\_percolated\_cluster\_meshes} + \texttt{plot\_percolation\_windows} & \texttt{Figure3JS.plot\_percolated\_cluster\_meshes} & \texttt{show\_percolation\_threejs}\\
\end{tabular}
\end{center}

Plotly-parity trace names: \texttt{plot\_surface}, \texttt{plot\_cell}, \texttt{draw\_framework} / \texttt{plot\_framework}, \texttt{plot\_maxima}, \texttt{plot\_connections}, \texttt{plot\_windows}.

\textbf{\textbf{Working directory:}} mofography repo root (so \texttt{examples/TestCIFs/...} resolves).

\textbf{\textbf{Environment:}} \texttt{mofography} installed (\texttt{pip install -e .} or dev env); Rust extension built if using \texttt{dgrid\_from\_atoms\_rust}.
\section{Prerequisites}
\label{sec:orgc61c842}

\begin{itemize}
\item CIF under \texttt{examples/TestCIFs/} (default below: \texttt{FOTNIN\_clean.cif}).
\item \texttt{ase}, \texttt{plotly}, \texttt{numpy}, \texttt{mofography}.
\item Emacs: \texttt{C-c C-c} on each block (same \texttt{:session mofography-viz} keeps variables).
\end{itemize}
\section{Setup — imports and paths}
\label{sec:org33767c0}

\begin{verbatim}
from pathlib import Path

import mofography as mgr
from ase.io import read

from mofography.threejs_viz import Figure3JS
from mofography.workflow import (
    show_mof_plotly,
    show_mof_threejs,
    show_percolation_threejs,
)
from mofography.winding_percolation import (
    compute_percolation_history_winding,
    get_pore_dia_mgr_winding,
)

REPO = Path(".").resolve()  # run from repo root
OUT = REPO / "examples" / "viz_outputs"
OUT.mkdir(parents=True, exist_ok=True)

CIF = REPO / "examples" / "TestCIFs" / "FOTNIN_clean.cif"
SPACING = 0.25
MASK, H, THRESH = 1.5, 0.2, 1.5

print("CIF:", CIF)
print("OUT:", OUT)
\end{verbatim}
\section{Segmentation}
\label{sec:orga87f95e}

Same pipeline as the interactive app before plotting.

\begin{verbatim}
atoms = read(str(CIF))
fw_atoms = mgr.update_with_covalent_bonds(atoms.copy())

atoms = mgr.update_with_covalent_bonds(atoms)
dgrid = mgr.dgrid_from_atoms_rust(atoms, spacing=SPACING)
seg = mgr.do_periodic_segmentation_rust(
    dgrid,
    atoms,
    spacing=SPACING,
    mask_thickness=MASK,
    h=H,
    threshold_abs=THRESH,
    return_dgrid=True,
)

stem = CIF.stem
print("seg keys:", sorted(seg.keys()))
print("dgrid shape:", seg["dgrid"].shape)
\end{verbatim}
\section{Main visualization (interactive app tab)}
\label{sec:org0836850}

\subsection{Plotly — one-liner}
\label{sec:org816b5a4}

\begin{verbatim}
fig = show_mof_plotly(seg, atoms, framework_atoms=fw_atoms, stylize=True)
out_plotly = OUT / f"{stem}_main_plotly.html"
fig.write_html(str(out_plotly))
print("Wrote:", out_plotly)
\end{verbatim}
\subsection{Three.js — one-liner}
\label{sec:orgf374806}

\begin{verbatim}
show_mof_threejs(
    seg,
    atoms,
    framework_atoms=fw_atoms,
    output_html=OUT / f"{stem}_main_threejs.html",
    plot_framework_representation="mesh",  # or "instanced" / "scatter"
    plot_surface=True,
    plot_surface_pbc=True,
    plot_connections=True,
    plot_windows=True,
    plot_maxima=True,
    plot_layout="dark",
)
print("Wrote:", OUT / f"{stem}_main_threejs.html")
\end{verbatim}
\subsection{Three.js — granular \texttt{Figure3JS} chain}
\label{sec:orgd81109f}

Explicit control; each \texttt{plot\_*} adds legend + opacity slider in the HTML HUD.

\begin{verbatim}
(
    Figure3JS()
    .plot_framework(fw_atoms, representation="mesh")
    .plot_cell(atoms, color="purple")
    .plot_surface(
        seg["dgrid"],
        seg["regions"],
        atoms,
        isovalue=1.0,
        name="Regions",
        opacity=1.0,
    )
    .plot_surface(
        seg["dgrid"],
        seg.get("regions_pbc"),
        atoms,
        isovalue=1.0,
        name="Regions (PBC)",
        opacity=1.0,
        trace_id="surface_pbc",
    )
    .plot_connections(
        seg["connections"],
        seg["all_maxima"],
        seg["dgrid"].shape,
        atoms,
        hide_periodic=True,
    )
    .plot_windows(seg, atoms=atoms)
    .plot_maxima(seg, atoms)
    .set_camera(eye=(2.0, 0.0, 1.35))
    .to_html(OUT / f"{stem}_main_granular_threejs.html")
)
print("Wrote:", OUT / f"{stem}_main_granular_threejs.html")
\end{verbatim}
\section{PLD and percolation analysis}
\label{sec:orgd5fa988}

Build the \texttt{analysis} dict (same as app \texttt{\_compute\_pld\_analysis}).

\begin{verbatim}
dgrid = seg["dgrid"]
lcd, pld, dfsd, _ = get_pore_dia_mgr_winding(
    mat_atoms=atoms,
    dgrid=dgrid,
    kwargs=dict(
        spacing=SPACING,
        mask_thickness=MASK,
        h=H,
        threshold_abs=THRESH,
        min_distance=1.5,
        extra_padding=5.0,
    ),
    results=seg,
)
percolation_history, pld_locations, _final = compute_percolation_history_winding(
    seg, step=0.05, expand_components=False
)

unique_regions = seg.get("unique_regions", [])
connections = [
    c for c in seg.get("connections", [])
    if all(x in unique_regions for x in c[0])
]
maxima_raw = seg.get("all_maxima", [])
regions = seg["regions"]
maxima = maxima_raw[: int(regions.max())] if len(maxima_raw) else maxima_raw
maxima_info = seg.get("maxima_info")
maxima_radii = (
    dgrid[tuple(maxima_info[:, 1:4].T)] if maxima_info is not None else None
)
rag = mgr.rag_from_connections(connections, maxima)
if maxima_radii is not None:
    rag = mgr.add_maxima_to_rag(rag, maxima, maxima_radii, regions.shape, atoms)
rag.graph["regions"] = regions
rag.graph["ase_atoms"] = atoms

analysis = {
    "lcd": float(lcd),
    "pld": float(pld),
    "dfsd": float(dfsd),
    "percolation_history": percolation_history,
    "pld_locations": pld_locations,
    "dgrid": dgrid,
    "regions": regions,
    "connections": connections,
    "maxima": maxima,
    "maxima_info": maxima_info,
    "maxima_pbc_dict": seg.get("maxima_pbc_dict", {}),
    "rag": rag,
    "isovalue": 1.0,
}

probe_r = analysis["percolation_history"][-1][0]
print(f"LCD={lcd:.3f} PLD={pld:.3f} probe_r={probe_r:.3f}")
\end{verbatim}
\section{PLD / percolation visualization}
\label{sec:org1eb80b9}

\subsection{Three.js — one-liner}
\label{sec:orgd065abf}

\begin{verbatim}
show_percolation_threejs(
    analysis,
    atoms,
    probe_radius=float(probe_r),
    output_html=OUT / f"{stem}_pld_percolation_threejs.html",
    layout="dark",
)
print("Wrote:", OUT / f"{stem}_pld_percolation_threejs.html")
\end{verbatim}
\subsection{Three.js — granular \texttt{Figure3JS} chain}
\label{sec:orgd3ce56d}

\begin{verbatim}
(
    Figure3JS()
    .plot_cell(atoms)
    .plot_connections(
        analysis["connections"],
        analysis["maxima"],
        analysis["dgrid"].shape,
        atoms,
        dgrid=analysis["dgrid"],
        regions=analysis["regions"],
    )
    .plot_maxima(
        analysis["dgrid"],
        analysis["maxima"],
        atoms,
        maxima_info=analysis.get("maxima_info"),
    )
    .plot_percolated_cluster_meshes(
        analysis["percolation_history"],
        analysis["dgrid"],
        analysis["regions"],
        atoms,
        probe_radius=float(probe_r),
        isovalue=analysis["isovalue"],
    )
    .plot_percolation_windows(analysis, atoms, probe_radius=float(probe_r))
    .to_html(OUT / f"{stem}_pld_granular_threejs.html")
)
print("Wrote:", OUT / f"{stem}_pld_granular_threejs.html")
\end{verbatim}
\subsection{Plotly — granular traces (optional)}
\label{sec:org60bb030}

If you prefer Plotly for the PLD tab, assemble traces like the app Plotly branch:

\begin{verbatim}
import plotly.graph_objects as go
import plotly.express as px
from mofography.visualize import (
    dark_layout,
    plot_cell,
    plot_maxima,
    plot_percolated_cluster_meshes,
    plot_percolation_windows,
    plot_rag,
)

traces = []
traces.append(plot_cell(atoms))
traces.extend(plot_rag(analysis["rag"], name="Pore Graph"))
traces.extend(
    plot_maxima(
        analysis["dgrid"],
        analysis["maxima"],
        atoms,
        maxima_info=analysis.get("maxima_info"),
        maxima_pbc_dict=analysis.get("maxima_pbc_dict"),
    )
)
traces.extend(
    plot_percolated_cluster_meshes(
        analysis["percolation_history"],
        analysis["dgrid"],
        analysis["regions"],
        atoms,
        probe_radius=float(probe_r),
        isovalue=analysis["isovalue"],
        stepsize=2,
        cscale=px.colors.qualitative.Pastel,
    )
)
traces.extend(
    plot_percolation_windows(analysis, probe_radius=float(probe_r), pld_tolerance=0.25)
)

fig_pld = go.Figure(data=traces)
fig_pld.update_layout(**dark_layout)
out_pld_plotly = OUT / f"{stem}_pld_plotly.html"
fig_pld.write_html(str(out_pld_plotly))
print("Wrote:", out_pld_plotly)
\end{verbatim}
\section{Trace API quick reference}
\label{sec:orgdb36d87}

\begin{center}
\begin{tabular}{llll}
App layer & Plotly (\texttt{visualize.py}) & Three.js functional & Three.js chain (\texttt{Figure3JS})\\
\hline
Framework mesh & \texttt{plot\_framework\_as\_mesh} & \texttt{draw\_framework(..., representation}``mesh'')= & \texttt{.plot\_framework(..., representation}``mesh'')=\\
Framework scatter & \texttt{draw\_framework} & \texttt{draw\_framework(..., representation}``scatter'')= & \texttt{.plot\_framework(..., representation}``instanced'')=\\
Unit cell & \texttt{plot\_cell} & \texttt{plot\_cell} & \texttt{.plot\_cell}\\
Distance surface & \texttt{plot\_surface\_dgrid} & \texttt{plot\_surface\_dgrid} & \texttt{.plot\_surface}\\
Maxima & \texttt{plot\_maxima} & \texttt{plot\_maxima} & \texttt{.plot\_maxima}\\
Pore graph & \texttt{plot\_connections\_as\_it\_is} & \texttt{plot\_connections} & \texttt{.plot\_connections}\\
Windows & \texttt{plot\_windows} & \texttt{plot\_windows} & \texttt{.plot\_windows}\\
Percolation clusters & \texttt{plot\_percolated\_cluster\_meshes} & \texttt{plot\_percolated\_cluster\_meshes} & \texttt{.plot\_percolated\_cluster\_meshes}\\
Percolation windows & \texttt{plot\_percolation\_windows} & \texttt{plot\_percolation\_windows} & \texttt{.plot\_percolation\_windows}\\
\end{tabular}
\end{center}
\section{Notes}
\label{sec:org4ede7ee}

\begin{itemize}
\item Change \texttt{CIF} in the Setup block to any structure under \texttt{examples/TestCIFs/}.
\item HTML files land in \texttt{examples/viz\_outputs/}; open in a browser (double-click or \texttt{open} on macOS).
\item To reset the Python session: \texttt{M-x org-babel-kill-session RET mofography-viz RET}
\item Internal trace IDs (e.g. \texttt{percolation\_clusters\_1\_x}) may be renamed in a future API cleanup; Plotly-parity method names stay stable.
\end{itemize}
\end{document}
